\section{Related work}
\label{sec:related}

\subsection{Public chest radiograph collections and deep classifiers}
Large public collections made deep learning on chest radiographs possible. NIH ChestX-ray14 released over one hundred thousand frontal radiographs with fourteen findings mined from radiology reports by natural-language processing~\cite{wang2017chestxray8}. CheXNet fine-tuned a DenseNet-121 on this collection~\cite{rajpurkar2017chexnet}, and CheXpert and MIMIC-CXR added larger hospital collections with report-derived labels and explicit uncertainty handling~\cite{irvin2019chexpert,johnson2019mimic}. VinDr-CXR provides radiologist annotations from several readers per image~\cite{nguyen2022vindr}, and the JSRT database contains nodule cases confirmed by computed tomography~\cite{shiraishi2000development}. For TB, the Montgomery County and Shenzhen sets~\cite{jaeger2014two,candemir2014lung} and TBX11K~\cite{liu2020rethinking} are the most common public sources, and the paediatric pneumonia set of Kermany et al.~\cite{kermany2018identifying} is widely used for pneumonia. Each collection differs in population, acquisition hardware, image processing and label provenance, and these differences are the starting point of the confounding problem studied here. Report-mined labels are known to be noisy; Oakden-Rayner documented label quality problems in large public collections and argued that they limit what models trained on them can show~\cite{oakdenrayner2020exploring}.

\subsection{COVID-19 radiograph classifiers and their critiques}
Many COVID-19 classifiers were trained on pooled collections in which COVID-19 images came from different sources than the other classes~\cite{wang2020covidnet,apostolopoulos2020covid,chowdhury2020can,rahman2021exploring}. The COVID-19 image data collection of Cohen et al.\ gathered images from publications and websites~\cite{cohen2020covid}, and the COVID-19 Radiography Database aggregated several such collections~\cite{chowdhury2020can,rahman2021exploring}. Critiques followed. Maguolo and Nanni showed that classifiers could separate the source datasets even when the lung region was removed from the images~\cite{maguolo2020critic}. DeGrave et al.\ showed that COVID-19 detectors relied on confounding factors rather than pathology and were likely to fail at new hospitals~\cite{degrave2021ai}. Two systematic reviews concluded that most published COVID-19 prediction models had a high risk of bias, often from data leakage, unclear provenance or the lack of external validation~\cite{wynants2020prediction,roberts2021common}. Related studies on chest radiograph classification for COVID-19, TB and pneumonia have appeared in this journal~\cite{nayak2021application,iqbal2023tuberculosis,amin2024automated,arias2024analysis}. Our COVID-19 class is subject to exactly these risks, which is why we report it per source and test the model's behaviour on an external cohort that contains no COVID-19 cases.

\subsection{Generalisation, shortcuts and bias}
Zech et al.\ showed that pneumonia classifiers trained on radiographs from one hospital system performed worse on other systems, and that the networks could identify the originating hospital from the image~\cite{zech2018variable}. Geirhos et al.\ framed such behaviour as shortcut learning: a model finds a decision rule that works on the test distribution but not on the intended task~\cite{geirhos2020shortcut}. Seyyed-Kalantari et al.\ showed that chest radiograph classifiers trained on public collections under-diagnose some patient subgroups~\cite{seyyedkalantari2021underdiagnosis}. Oakden-Rayner et al.\ showed that aggregate metrics can hide poor performance on clinically important subsets of a class (hidden stratification)~\cite{oakdenrayner2020hidden}, and a systematic review of deep learning algorithms in radiology found that most showed lower performance on external data than on internal data~\cite{yu2022external}. Domain shift and cross-domain generalisation in chest radiograph classification have been studied~\cite{pooch2020can,cohen2020limits}, and dataset bias has been analysed for COVID-19 classifiers trained on small pooled collections~\cite{tartaglione2020unveiling}. Beyond medical imaging, Kapoor and Narayanan surveyed data leakage in machine-learning-based science and proposed a taxonomy of its forms~\cite{kapoor2023leakage}. These studies motivate three design choices of the present work: grouping by patient and by near-duplicate content, reporting results per source, and evaluating on a hospital whose images were never used for training or selection.

\subsection{Contrast enhancement and filtering}
Adaptive histogram equalisation and its contrast-limited variant (CLAHE) are standard enhancement steps for radiographs~\cite{pizer1987adaptive,zuiderveld1994contrast}. Rahman et al.\ compared several enhancement methods, including CLAHE and gamma correction, for COVID-19 radiograph classification~\cite{rahman2021exploring}. Our preprocessing combines a light Gaussian pre-blur, CLAHE and a Gaussian low-pass filter. We describe it precisely (Section~\ref{sec:preproc}) but do not claim that it improves generalisation; a comparison with no filtering on the rebuilt data is reported in Section~\ref{sec:res-ablation}.

\subsection{Reporting, calibration and uncertainty}
The CLAIM checklist and TRIPOD+AI ask authors to report data sources, inclusion and exclusion rules, partitioning, the handling of patients with several images, external validation and calibration~\cite{mongan2020claim,collins2024tripodai}. Modern networks are often over-confident; the expected calibration error (ECE) summarises the gap between confidence and accuracy~\cite{naeini2015obtaining,guo2017calibration}. For proportions estimated on small cohorts we use the Wilson score interval~\cite{wilson1927probable}, and for composite metrics the percentile bootstrap~\cite{efron1994introduction}. Gradient-weighted class activation mapping (Grad-CAM) gives a coarse localisation of the image regions that influence a class score~\cite{selvaraju2017gradcam}; we use it only qualitatively.

\subsection{Position of this work}
The present study does not introduce a new learning method, and we do not compare our numbers with those of other papers, because they were obtained on different test sets. Its contribution is a documented, reproducible construction of a multi-source six-class benchmark with leakage controls, and a measurement of the gap between internal and external behaviour of a standard transfer-learning classifier trained on it.
